\documentclass[12pt]{article}
\usepackage{amsmath,amssymb,amsfonts,amsthm}
\usepackage[margin=1in]{geometry}

\begin{document}

\begin{center}
{\Large\bfseries Chapter 4, Problem 14}\\[6pt]
Byron \& Fuller, \textit{Mathematics of Classical and Quantum Physics}
\end{center}

\bigskip

\textbf{Problem.} Derive the matrix $\Lambda$ of the general Lorentz transformation corresponding to relative velocity $v$ in a direction given by the spherical angles $\theta, \phi$ by a similarity transformation, $\Lambda = R^{-1}\Lambda'R$, where $R$ is the rotation matrix, $R(\theta, \phi)$, of Eq.\ (1.23) on a 4-element that acts only on the three space coordinates. Note that $R^\dagger = R$. Assume that the axes of $K$ and $K'$ remain parallel during their uniform relative motion.

\bigskip
\textbf{Solution.}

The standard Lorentz boost along the $x$-axis is:
\[
\Lambda' = \begin{pmatrix} \gamma & -\beta\gamma & 0 & 0 \\ -\beta\gamma & \gamma & 0 & 0 \\ 0 & 0 & 1 & 0 \\ 0 & 0 & 0 & 1 \end{pmatrix},
\]
where $\beta = v/c$ and $\gamma = (1-\beta^2)^{-1/2}$.

To boost in a general direction specified by spherical angles $(\theta, \phi)$, we first rotate coordinates so that the boost direction aligns with the $x$-axis, apply the standard boost, then rotate back.

The rotation matrix $R(\theta,\phi)$ acts on the spatial coordinates and takes the direction $(\theta,\phi)$ to the $x$-axis. In the 4-dimensional spacetime, it takes the form:
\[
\mathcal{R} = \begin{pmatrix} 1 & 0 \\ 0 & R(\theta,\phi) \end{pmatrix},
\]
where $R(\theta,\phi)$ is the $3\times 3$ rotation matrix from Eq.\ (1.23).

The general Lorentz boost is then:
\[
\Lambda = \mathcal{R}^{-1}\,\Lambda'\,\mathcal{R}.
\]

Since $R$ is a real orthogonal matrix, $R^{-1} = R^T$. If the velocity direction has unit vector $\hat{n} = (\sin\theta\cos\phi,\, \sin\theta\sin\phi,\, \cos\theta)$, then the resulting general boost matrix has components:
\[
\Lambda^{0}{}_{0} = \gamma, \quad \Lambda^{0}{}_{i} = \Lambda^{i}{}_{0} = -\beta\gamma\, n_i,
\]
\[
\Lambda^{i}{}_{j} = \delta_{ij} + (\gamma - 1)\,n_i n_j,
\]
where $n_i$ are the components of $\hat{n}$.

This can be verified by direct computation. The spatial part is:
\[
\Lambda^{i}{}_{j} = (R^{-1})^{i}{}_{k}\,\Lambda'^{k}{}_{l}\,R^{l}{}_{j}.
\]
The only nontrivial spatial element of $\Lambda'$ is $\Lambda'^{1}{}_{1} = \gamma$; all other diagonal spatial elements are 1. Writing $R^{1}{}_{j} = n_j$ (the first row of $R$ projects onto $\hat{n}$):
\[
\Lambda^{i}{}_{j} = \delta_{ij} + (\gamma-1)\,n_i\,n_j,
\]
which is the standard result for a general Lorentz boost. \qed

\end{document}
